% Appendix H: Admissibility Functionals and Constraint Closure

\section*{H.1 Introduction}

Projection alone generates possibilities. A second structure is required
to evaluate them. This structure is called \emph{admissibility}.

\begin{tcolorbox}[enhanced, colback=RSVPdeep, colframe=RSVPdeep,
  coltext=white, arc=6pt, boxrule=0pt, left=6mm, right=6mm, top=5mm, bottom=5mm]
  Projection answers: \emph{What can be produced?}\\
  Admissibility answers: \emph{What should be accepted?}
\end{tcolorbox}

\begin{definition}{Admissibility Functional}{adm-func}
  An \emph{admissibility functional} is a map $A : X \to \mathbb{R}$.
  Large values indicate coherence; small values indicate violation.
  The \emph{admissible region} at threshold $\tau$ is
  $\mathcal{A}_\tau = \{x : A(x) \geq \tau\}$.
\end{definition}

\begin{example}
  With constraint violation $\kappa$, set $A(x) = e^{-\kappa(x)}$.
  Then $\kappa = 0 \Rightarrow A = 1$; large violations imply $A \to 0$.
\end{example}

\begin{definition}{Local and Global Constraints}{local-global}
  A \emph{local constraint} depends only on a neighborhood: $\kappa_L(x)$.
  A \emph{global constraint} depends on an entire trajectory: $\kappa_G(\gamma)$.
  Admissibility typically combines both: $A = A_L \cdot A_G$.
\end{definition}

\begin{definition}{Closure Operator}{closure-op}
  A \emph{closure operator} is $\mathrm{Cl} : \mathcal{P}(X) \to \mathcal{P}(X)$
  satisfying extensivity, idempotence, and monotonicity.
  Closure adds whatever structure is required for coherence.
\end{definition}

\begin{definition}{Obstruction}{obstruction-adm}
  An \emph{obstruction} is a quantity $\epsilon(x)$ measuring failure of
  closure (missing assumptions, contradictions, broken references).
  $A(x) = f(\epsilon(x))$; large obstructions imply low admissibility.
\end{definition}

\begin{definition}{Trajectory Admissibility}{traj-adm}
  For $\gamma : [0,T] \to X$, define
  $A(\gamma) = \int_0^T a(\gamma(t))\,dt$.
\end{definition}

\begin{definition}{Admissible Accessibility}{adm-access}
  $\mathcal{N}(x) = \{y : A(y) \geq \tau\}$, with
  $S_A(x) = \log |\mathcal{N}(x)|$.
  This quantity becomes central in RSVP (Appendix~K).
\end{definition}

\begin{theorem}{Projection Does Not Determine Validity}{proj-not-valid}
  Projection alone cannot determine validity. A second evaluation
  structure — admissibility — is required.
\end{theorem}
\begin{proof}
  There exist outputs satisfying projection while violating external
  constraints; admissibility is the only mechanism to filter them.
\end{proof}

\begin{namedthm}{The Admissibility Principle}
  Every intelligent system requires two distinct layers:
  \emph{projection} $X \to Y$ (generating possibilities) and
  \emph{admissibility} $A : Y \to \mathbb{R}$ (evaluating them).
  Generation without admissibility produces chaos.
  Admissibility without generation produces stagnation.
  Intelligence emerges from their interaction.
\end{namedthm}
